\documentclass[11pt]{article}

\usepackage[margin=1in]{geometry}
\usepackage{amsmath, amssymb}
\usepackage{booktabs}
\usepackage{hyperref}
\usepackage{url}
\usepackage{enumitem}
\setlist{leftmargin=*}
\setlength{\parskip}{0.35em}
\setlength{\parindent}{0pt}

\title{\textbf{Stage 1: Background and Related Work}\\
Scalable Graph Partitioning for Communication-Efficient Distributed Optimization}
\author{Team Members: Name 1, Name 2, Name 3}
\date{CS240: Algorithm Design and Analysis, Spring 2026}

\begin{document}

\maketitle

\section{Course-Level Algorithmic Background}

Before introducing the application, it is useful to recall several basic ideas
from algorithm design and analysis. Many modern data problems can be made
precise by first choosing a mathematical representation. In this project, the
representation is a graph. A graph abstracts a system into vertices and edges:
vertices are objects or computational units, while edges represent
relationships, dependencies, or communication links. If edges have weights, the
weights usually represent cost, similarity, capacity, or strength of
interaction.

Graph problems are central in algorithms because they connect modeling,
optimization, and complexity. Shortest paths, minimum spanning trees, matchings,
flows, cuts, and graph partitioning are all examples where a concrete
application becomes an algorithmic problem after the graph model is chosen.
Among these, cut problems are especially relevant here. A cut separates a graph
into different vertex sets and measures the edges crossing between them.
Minimum-cut algorithms can solve some cut problems exactly, but adding extra
requirements such as balance makes the problem substantially harder.

This project therefore starts from a familiar course idea---representing a
problem as a graph---and then studies a harder version of a cut problem:
balanced graph partitioning. The important algorithmic themes are:
\begin{itemize}
    \item \textbf{objective functions}: minimizing cross-partition edge weight;
    \item \textbf{constraints}: keeping partitions balanced so that no worker
    receives too much work;
    \item \textbf{heuristics}: using practical methods when exact optimization
    is too expensive;
    \item \textbf{complexity and scalability}: studying how runtime changes as
    the graph becomes larger.
\end{itemize}

With this background, the project can be read as a step-by-step application of
course material. Stage 1 builds the graph model and explains the related work.
Stage 2 chooses and analyzes a representative algorithm. Later stages implement
the algorithm, test it, and discuss possible extensions.

\section{Application Background}

This project studies how classical graph partitioning algorithms can be used in
communication-efficient distributed optimization. In modern machine learning,
scientific computing, and large-scale data processing, a problem is often too
large to be solved on a single machine. A common solution is to distribute the
work across several workers or agents. Each worker owns part of the data, part
of the variables, or part of the computation, and the workers exchange messages
to coordinate their local updates.

Communication is often one of the main bottlenecks in such systems. Local
floating-point computation can be cheap compared with synchronization,
cross-machine data movement, and repeated message passing. If two strongly
dependent tasks are assigned to different workers, their interaction becomes a
cross-worker communication event. If many such edges cross worker boundaries,
the distributed algorithm may spend substantial time moving information rather
than computing.

Graph partitioning provides a natural algorithmic abstraction for this issue.
We model the computation as a graph $G=(V,E,w)$. Vertices represent data
blocks, variable blocks, computational tasks, or agents. Edges represent
dependencies, message-passing requirements, or communication links. Edge
weights encode the strength or cost of communication. A partition of $V$ into
$k$ balanced parts corresponds to assigning tasks to $k$ workers. The desired
partition should keep the workload balanced while minimizing the total weight
of edges crossing between parts.

\section{Algorithmic Formulation}

The main computational problem is balanced $k$-way graph partitioning. Given an
undirected weighted graph $G=(V,E,w)$ and an integer $k$, find disjoint sets
$V_1,\ldots,V_k$ such that
\[
    V = V_1 \cup \cdots \cup V_k,
    \qquad V_i \cap V_j = \emptyset \text{ for } i\neq j.
\]
The cut objective is
\[
    \mathrm{cut}(V_1,\ldots,V_k)
    =
    \sum_{(u,v)\in E,\; u\in V_i,\; v\in V_j,\; i\neq j} w_{uv}.
\]
To prevent one worker from receiving almost all vertices, we also require a
balance condition such as
\[
    |V_i| \leq (1+\epsilon)\frac{|V|}{k},
    \qquad i=1,\ldots,k.
\]
The optimization goal is to minimize the cut subject to this balance
constraint.

This formulation connects directly to classical algorithms. Minimum cut and
maximum flow solve some exact cut problems, but balanced graph partitioning is
harder because the balance constraint prevents trivial solutions. Spectral
methods use eigenvectors of the graph Laplacian to find low-conductance
separations. Local-search methods such as Kernighan--Lin iteratively improve a
partition by swapping vertices. Multilevel methods such as METIS improve
scalability by coarsening the graph, partitioning a smaller graph, and refining
the result on the original graph.

\section{Representative Literature}

\textbf{Kernighan and Lin.}
Kernighan and Lin proposed one of the classical heuristic procedures for graph
partitioning. Their method starts from an initial bisection and repeatedly
selects vertex swaps that reduce the cut. Although the original problem is
graph bisection rather than general $k$-way partitioning, the core local-search
idea remains highly influential. Our project reproduces this refinement idea
and adapts it to a practical $k$-way setting by applying pairwise
Kernighan--Lin bisection steps between parts.

\textbf{Spectral graph partitioning.}
Spectral partitioning uses the graph Laplacian $L=D-A$, where $D$ is the degree
matrix and $A$ is the adjacency matrix. The Fiedler vector, the eigenvector
corresponding to the second-smallest eigenvalue of $L$, provides a relaxed
solution to a graph cut problem. Sorting vertices by the Fiedler vector and
splitting them gives a natural bisection. Recursive bisection then produces
multiple parts. This method is important because it reveals a deep connection
between graph cuts, linear algebra, and continuous relaxations.

\textbf{Multilevel partitioning.}
Karypis and Kumar's METIS-style multilevel framework improved practical graph
partitioning quality and speed. The graph is first coarsened into a sequence of
smaller graphs. A partition is computed on the coarsest graph, then projected
back and refined. We do not implement the full multilevel pipeline in the core
project, but METIS is a key reference point because it shows how local
refinement can be made scalable for large irregular graphs.

\textbf{Distributed optimization over networks.}
Distributed subgradient, consensus, and decentralized gradient methods show
how optimization algorithms behave when computation is spread across agents
connected by a communication topology. In these methods, the graph controls
which agents exchange information. Our project uses this setting as the modern
application domain: graph partitioning is not only a graph-processing task, but
also a tool for reducing cross-worker communication in distributed algorithms.

\section{Why This Topic Is Algorithmically Interesting}

This topic is algorithmically interesting for four reasons. First, balanced
graph partitioning combines an objective function with a hard feasibility
constraint. A partition with tiny cut but terrible balance is not useful for
distributed computation, while a perfectly balanced partition with many
cross-edges may communicate too much. The algorithm must trade off both goals.

Second, the problem illustrates the difference between exact algorithms,
relaxations, and heuristics. The project compares a greedy baseline, a spectral
relaxation, and a local-search refinement method. These methods represent
different algorithm design strategies taught in an algorithms course.

Third, the application is modern. Graph partitioning appears in distributed
machine learning, graph neural network training, sparse matrix computations,
and multi-agent optimization. A classical algorithmic tool therefore helps
solve a current systems problem.

Finally, the project is experimentally measurable. We can evaluate cut weight,
normalized cut, balance ratio, runtime, scalability, and cross-partition
communication in a simple consensus simulation. These metrics connect the
formal objective to observable system behavior.

\section{Project Positioning}

Our selected reference topic is \textbf{Topic 2: Community Detection / Graph
Partitioning}. Instead of focusing on social communities, we focus on graph
partitioning for distributed optimization. This choice still fits the reference
topic because the central algorithmic problem is graph partitioning, but it
connects the problem to a more modern data and machine learning application.

The main method selected for reproduction is the Kernighan--Lin graph
partitioning heuristic. Spectral bisection and greedy balanced partitioning are
implemented as baselines. The expected outcome of the project is to demonstrate
that classical graph algorithms can reduce cross-worker communication while
maintaining reasonable load balance.

\begin{thebibliography}{9}

\bibitem{kernighan1970}
B. W. Kernighan and S. Lin.
\newblock An efficient heuristic procedure for partitioning graphs.
\newblock \emph{Bell System Technical Journal}, 1970.

\bibitem{karypis1998}
G. Karypis and V. Kumar.
\newblock A fast and high quality multilevel scheme for partitioning irregular graphs.
\newblock \emph{SIAM Journal on Scientific Computing}, 1998.

\bibitem{nedic2009}
A. Nedic and A. Ozdaglar.
\newblock Distributed subgradient methods for multi-agent optimization.
\newblock \emph{IEEE Transactions on Automatic Control}, 2009.

\end{thebibliography}

\end{document}
